% Section 1.3, from lectures/L01/README.md: the objectives, the questions,
% the further reading, and what the next chapter does with all of it.


\section{Review}
\label{c1:sec:review}

\needspace{6\baselineskip}
\subsection*{What you should be able to do now}
After this chapter, you should be able to:
\begin{itemize}
  \item State the sign convention this book uses, and say what breaks if it is inverted.
  \item Compute a divider's output, its Thevenin resistance, and its output under any load, by hand.
  \item Say why a 10 kilohm load on a 33k/6.8k divider is not a light load.
  \item Write the node equation for any node in a resistive network without deriving it from scratch.
  \item Stamp a resistor, a current source and a voltage source into a matrix.
  \item Explain why a voltage source needs an extra unknown, and what that unknown is.
  \item Implement \code{ael::net} and \code{ael::mna} against the shipped suite.
\end{itemize}

\needspace{6\baselineskip}
\subsection*{Questions to test yourself}
You should be able to answer:
\begin{enumerate}
  \item A divider made of 1 megohm and 1 megohm, and one made of 10 ohms and 10 ohms, have the same ratio. What is different about them, and which would you put on an oscilloscope probe?
  \item Why is a divider's Thevenin resistance the same whichever leg you call the upper one?
  \item Nodal analysis gives one equation per node. A network has five nodes. How many unknowns are there, and why is it not five?
  \item You stamp a resistor between two nodes and get four matrix entries. Two are positive and two are negative. Which are which, and what would swapping them mean physically?
  \item A voltage source between two nodes cannot be written as a conductance. Why not, and what does the extra row in the matrix say?
  \item Your solver returns node voltages that are all zero for a network you know is live. Name two causes, and say which one the shipped suite catches.
  \item Superposition says two sources add. Your solver never mentions superposition. Why does it hold anyway, and in which chapter does it stop holding?
\end{enumerate}

\needspace{6\baselineskip}
\subsection*{Further reading}
\begin{itemize}
  \item \emph{The Art of Electronics}, Horowitz and Hill, chapter 1, for a second treatment of the same material at roughly the same level.
\end{itemize}

\needspace{6\baselineskip}
\subsection*{Looking ahead}
\begin{itemize}
  \item Chapter~\ref{c2:ch:reactance} makes the solver complex, which is the whole of the frequency domain and costs about thirty lines. Reactance, phasors, and the Bode plot as two straight lines and a corner.
\end{itemize}
